\section{A Mutation That Can Fail}
\label{sec:ch04-a-mutation-that-can-fail}

\index{mutation!the book's first}%
A query that gains a field inconveniences nobody. A mutation gains a way to
fail every time somebody adds a business rule, which is why the book's first
one is here rather than in chapter~\ref{ch:hot-chocolate-zero}: where those
failures land in the response is a decision you make once and live with.

\index{rescheduling!the operation}%
The operation moves a session to a different slot, and it can fail two ways
the caller should be able to act on. The session may not exist. The new slot
may collide with another session the same speaker is giving. Both are answers
rather than accidents.

\index{exception!as a declared failure}%
Each is a C\# exception, and each exception is the shape of the error the
client will see, so its properties are chosen for the client rather than for
the log. The missing session first, in
\code{SessionNotFoundException.cs}:

\begin{minted}{csharp}
using HotChocolate.Types.Relay;

namespace Sessions;

/// <summary>
/// No session on the program carries the identifier that was asked for.
/// </summary>
public sealed class SessionNotFoundException(int sessionId)
    : Exception($"No session has id {sessionId}.")
{
    /// <summary>
    /// The identifier that matched nothing.
    /// </summary>
    [ID<Session>]
    public int SessionId { get; } = sessionId;
}
\end{minted}

Then \code{SpeakerDoubleBookedException.cs}:

\begin{minted}{csharp}
using HotChocolate.Types.Relay;

namespace Sessions;

/// <summary>
/// The new slot overlaps another session the same speaker gives.
/// </summary>
public sealed class SpeakerDoubleBookedException(
    int speakerId,
    int clashingSessionId)
    : Exception(
        $"Speaker {speakerId} already gives session "
        + $"{clashingSessionId} in that slot.")
{
    /// <summary>
    /// The speaker who cannot be in two rooms at once.
    /// </summary>
    [ID<Speaker>]
    public int SpeakerId { get; } = speakerId;

    /// <summary>
    /// The session already occupying the slot.
    /// </summary>
    [ID<Session>]
    public int ClashingSessionId { get; } = clashingSessionId;
}
\end{minted}

\index{IDattribute@\code{[ID<T>]}!on an error type}%
Those \code{[ID<T>]} attributes are the section that just ended applying to a
type it did not obviously reach. An error is part of the schema, so an error
carries node ids like everything else, and a client that receives
\enquote{session 4 is already in that slot} can pass the id straight back to
\code{node} and show the reader which session. Handing back a bare \code{4}
would also have republished the primary key that
chapter~\ref{ch:data-without-n-plus-1} took the trouble to hide.

\index{Mutation@\code{Mutation}!the type class}%
\code{Mutation.cs} is one method, and the two attributes above it are what
turn the exceptions into schema:

\begin{minted}{csharp}
using Microsoft.EntityFrameworkCore;

namespace Sessions;

[MutationType]
public static class Mutation
{
    /// <summary>
    /// Moves one session to a new slot, keeping its length.
    /// </summary>
    [Error<SessionNotFoundException>]
    [Error<SpeakerDoubleBookedException>]
    public static async Task<Session> RescheduleSessionAsync(
        [ID<Session>] int sessionId,
        DateTimeOffset startsAt,
        ConferenceContext db,
        CancellationToken ct)
    {
        var session = await db.Sessions
            .FirstOrDefaultAsync(s => s.Id == sessionId, ct)
            ?? throw new SessionNotFoundException(sessionId);

        var ends = startsAt.AddMinutes(session.DurationMinutes);

        // The overlap test is done in memory on purpose. StartsAt goes
        // through a value converter, so SQLite cannot do arithmetic on
        // that column, and one speaker's program is a handful of rows.
        var alsoTheirs = await db.Sessions
            .Where(s => s.SpeakerId == session.SpeakerId)
            .Where(s => s.Id != session.Id)
            .ToListAsync(ct);

        var clash = alsoTheirs.FirstOrDefault(
            s => s.StartsAt < ends
                && startsAt < s.StartsAt.AddMinutes(s.DurationMinutes));

        if (clash is not null)
        {
            throw new SpeakerDoubleBookedException(
                session.SpeakerId,
                clash.Id);
        }

        var moved = session with { StartsAt = startsAt };
        db.Entry(session).CurrentValues.SetValues(moved);
        await db.SaveChangesAsync(ct);

        return moved;
    }
}
\end{minted}

\index{value converter!blocks translation}%
The comment in the middle is the value converter from
chapter~\ref{ch:data-without-n-plus-1} coming back to collect. The start time
is stored as an instant and converted on the way out, so EF~Core cannot
translate arithmetic over that column into SQL, and an overlap test is
arithmetic. Loading one speaker's sessions and testing in memory is the honest
answer at this size, and it costs the second statement in the counts below.

\index{mutation conventions!what they generate}%
Now the part that is not in that file. \code{AddMutationConventions} on the
builder, one call, and the schema gains every one of these. As before, the
descriptions and the \code{@cost} directive the export writes are dropped,
and two definitions it writes on one line are wrapped to fit the measure:

\begin{graphqlsdl}
type Mutation {
  rescheduleSession(
    input: RescheduleSessionInput!
  ): RescheduleSessionPayload!
}

input RescheduleSessionInput {
  sessionId: ID!
  startsAt: DateTime!
}

type RescheduleSessionPayload {
  session: Session
  errors: [RescheduleSessionError!]
}

union RescheduleSessionError =
  SessionNotFoundError | SpeakerDoubleBookedError

interface Error {
  message: String!
}

type SessionNotFoundError implements Error {
  message: String!
  sessionId: ID!
}

type SpeakerDoubleBookedError implements Error {
  message: String!
  speakerId: ID!
  clashingSessionId: ID!
}
\end{graphqlsdl}

\index{mutation conventions!input type}%
The method's non-service parameters became the fields of an input type named
after the mutation, and the argument itself is called \code{input}. That is
not decoration: an argument list is positional in the schema's grammar and
adding to it is awkward, while an input object grows a field without any
existing query changing. The shape is Relay's, and Relay's current
documentation no longer mentions the \code{clientMutationId} field its older
editions asked every mutation to carry. Nothing generates one here.

\index{payload!success field is nullable}%
The payload's \code{session} field is nullable, and the method returns
\code{Task<Session>}, which in C\# with nullable reference types enabled is a
promise that it is never null. The convention makes the field nullable anyway,
and it is right to: the call can answer with errors instead, and a payload
that promised a session would be lying about two of its three outcomes. This
is the schema and the C\# annotation disagreeing on purpose, and
section~\ref{sec:ch04-what-a-non-null-field-promises} is where that turns into
a rule.

\index{error type!generated from an exception}%
Each \code{[Error<T>]} contributed an object type. The naming rule strips
\code{Exception} and appends \code{Error}, so the not-found exception became
\code{SessionNotFoundError}. That is the same rule, a generated name derived
from a name you wrote, that the two previous chapters met on
\code{AddSessionsTypes} and on the DataLoader interface. Inside the type, the
exception's message became a
\code{message} field. Every other public property became a field too, which is
the only reason the two id fields are in the schema at all.

\index{union!of error types}%
\index{abstract type!union versus interface}%
Then a union of those types, and an interface both of them implement. Either
one alone would be worse than the pair.

A union is a closed list. \code{errors} returns a
\code{RescheduleSessionError}, and the schema says exactly which types that can
be, so a client reading the schema knows every way this call can fail before
writing a line. A generated client can produce an exhaustive switch. That is
the expressiveness Marc-Andr\'{e} Giroux argues for in \emph{A Guide to GraphQL
Errors}~\autocite{giroux2020}: \enquote{Error unions are great because they are
a really expressive way of structuring our schema. It lets clients see right
away what could happen when querying or mutating a resource.}

The problem with a closed list is that adding to it changes what an existing
client's exhaustive switch covers. Which is where the interface earns its
place, and Giroux names that combination too:

\begin{quote}
What if we could combine the extensibility of the interface with the
expressivity of the union? Well, we can! [\ldots] Make sure all of your union
members implement a common interface for this solution to work.
\end{quote}

\index{Error@\code{Error} interface!fallback for unknown members}%
A client writes \code{... on Error \{ message \}} once, and every future member
of the union answers it, including the ones added after that client shipped.
The union carries what is knowable now, the interface carries what is knowable
forever, and a new failure mode becomes an additive change. Both fell out of
two attributes.

\index{errors as data!versus the errors array}%
An exception that is not declared escapes into the top-level \code{errors}
array instead, where it is a string with a path and no type, arriving beside
\code{data: null} for the whole field. A client can display it and cannot
branch on it, because there is nothing to branch on. Declaring the failure
moves it into the payload, where it has a type, has fields, and is something
the client can act on.

\subsection{The Three Outcomes}

\index{mutation!the operation the three runs send}%
All three runs below send this one operation and differ only in their
variables. Both error types are named in it so that the responses show their
own fields, which a client that only wanted to display the message would not
bother with.

\begin{minted}{graphql}
mutation R($input: RescheduleSessionInput!) {
  rescheduleSession(input: $input) {
    session { id title startsAt }
    errors {
      __typename
      ... on SessionNotFoundError { message sessionId }
      ... on SpeakerDoubleBookedError {
        message
        speakerId
        clashingSessionId
      }
    }
  }
}
\end{minted}

\index{mutation!a session that does not exist}%
On the wire it is an ordinary POST carrying that operation and a variables
object. The envelope for the first run, with the operation text standing in
for the block above and the body wrapped across three lines to fit the page:

\begin{minted}{text}
POST /graphql HTTP/1.1
Host: localhost:5001
Content-Type: application/json

{"query":"<the operation above>","variables":
{"input":{"sessionId":"U2Vzc2lvbjo5OQ==",
"startsAt":"2026-09-14T16:00:00Z"}}}
\end{minted}

That id decodes to \code{Session:99}, and there is no session 99. One
statement, the lookup, and then the exception:

\begin{minted}[fontsize=\footnotesize]{json}
{"data":{"rescheduleSession":{"session":null,"errors":[
{"__typename":"SessionNotFoundError","message":"No session has id 99.",
"sessionId":"U2Vzc2lvbjo5OQ=="}]}}}
\end{minted}

\index{errors as data!status 200}%
Status 200, no error at the top level, no null in \code{data}. As far as HTTP
is concerned nothing went wrong, because nothing did: the request was valid,
the server answered it, and the answer is that there is no such session.

\index{mutation!a clash}%
The second failure needs the schedule. Ada Fischer gives session~1 at nine and
session~2 at ten, forty-five minutes each, so moving session~2 to half past
nine puts her in two rooms at once.

\begin{minted}{text}
{"input":{"sessionId":"U2Vzc2lvbjoy",
          "startsAt":"2026-09-14T09:30:00Z"}}
\end{minted}

\begin{minted}[fontsize=\footnotesize]{json}
{"data":{"rescheduleSession":{"session":null,"errors":[
{"__typename":"SpeakerDoubleBookedError",
"message":"Speaker 1 already gives session 1 in that slot.",
"speakerId":"U3BlYWtlcjox","clashingSessionId":"U2Vzc2lvbjox"}]}}}
\end{minted}

\index{statement count!mutation}%
Two statements, and no third. The lookup, the speaker's other sessions, and
then nothing: the clash was detected before anything was written, so the
\code{UPDATE} never ran. That is the property that separates a declared failure
from an exception thrown halfway through a transaction, and it is a property of
how the method is written rather than anything the conventions did.

\index{mutation!success}%
A move that works costs three: the lookup, the speaker's other sessions, and
the write.

\begin{minted}[fontsize=\footnotesize]{json}
{"data":{"rescheduleSession":{"session":{"id":"U2Vzc2lvbjoz",
"title":"Paging Without Offsets","startsAt":"2026-09-14T16:00:00Z"},
"errors":null}}}
\end{minted}

\code{errors} is null rather than an empty list, so a client tests for absence
and not for length. The session comes back with the same node id it had before
the move, which is the identity work of
section~\ref{sec:ch04-every-entity-gets-a-name} paying off: rescheduling a
session does not make it a different session, and nothing about the id says
when it starts.

\begin{onhc14}
The mutation conventions, \code{[MutationType]}, \code{[Error<T>]} and
\code{[ID<T>]} on an exception property all exist on Hot Chocolate~14 and
generate the same input type, the same payload, the same union and the same
\code{Error} interface. Every file in this section compiles there unchanged,
and the three statement counts are the same.
Appendix~\ref{app:hc14} carries the full mapping.
\end{onhc14}
